% !TeX program = xelatex
% !TeX encoding = UTF-8
% academic-report.tex — 学术报告模板（多节结构）
% 适用于：组会报告、研讨会、课程展示、读书报告

\documentclass[
    aspectratio=169,
    fontset=windows,
]{ctexbeamer}

\mode<presentation>

\usetheme[maxplus,red,light,shadow]{sjtubeamer}

\usepackage{amsmath}
\usepackage{graphicx}
\usepackage{booktabs}
\usepackage{tikz}
\usepackage{fontawesome5}
\usepackage{multicol}

\graphicspath{{figures/}}

% 仅保留页码，去掉导航超链接
\setbeamertemplate{navigation symbols}{\insertframenumber/\inserttotalframenumber}

% 节目录
\AtBeginSection[]{
  \begin{frame}
    \tableofcontents[currentsection,subsectionstyle=show/show/hide]
  \end{frame}
}

% 标题 — 不建议用 \institute（会占用标题页空间导致溢出）
\title{学术报告标题}
%\subtitle{副标题（可选）}
\author{报告人姓名}
\date{\today}

\begin{document}

\maketitle

\begin{frame}{报告提纲}
  \tableofcontents[hideallsubsections]
\end{frame}

% ══════════════════════════════════════
\section{研究背景}
% ══════════════════════════════════════

\begin{frame}{研究背景}
  \begin{itemize}
    \item 研究领域的现状和重要性
    \item 当前存在的关键问题
    \item 本研究的出发点
  \end{itemize}
\end{frame}

% ══════════════════════════════════════
\section{研究方法}
% ══════════════════════════════════════

\begin{frame}{方法概述}
  \begin{block}{核心思路}
    用一句话概括方法的核心思想。
  \end{block}

  \begin{enumerate}
    \item 步骤一
    \item 步骤二
    \item 步骤三
  \end{enumerate}
\end{frame}

% ══════════════════════════════════════
\section{实验结果}
% ══════════════════════════════════════

\begin{frame}{实验设置}
  \begin{table}
    \centering
    \begin{tabular}{lcc}
      \toprule
      \textbf{方法} & \textbf{指标 A} & \textbf{指标 B} \\
      \midrule
      Baseline & 85.2 & 72.1 \\
      \textbf{本文方法} & \textbf{91.3} & \textbf{78.5} \\
      \bottomrule
    \end{tabular}
  \end{table}
\end{frame}

% ══════════════════════════════════════
\section{总结与展望}
% ══════════════════════════════════════

\begin{frame}{总结}
  \begin{enumerate}
    \item 贡献一
    \item 贡献二
    \item 贡献三
  \end{enumerate}
\end{frame}

\begin{frame}{未来工作}
  \begin{itemize}
    \item 方向一
    \item 方向二
    \item 方向三
  \end{itemize}
\end{frame}

% 尾页：重复封面（不使用 \makebottom，避免中英文混用问题）
\maketitle

\end{document}
